% Ad Familiares 7.8 --- headnote
% ad-familiares-07-08, end of July 54 BC, Rome

\textit{Cicero to C. Trebatius Testa, written from Rome at the
end of July 54 BC. The recipient is the young legal expert
whom Cicero had recommended to Caesar in early 54 (the famous
\textit{Fam.} 7.5, the recommendation letter, opens the Trebatius
sequence). Trebatius had been on Caesar's staff in Gaul through
the spring and summer of 54, on his way (with Caesar's army)
to the second British expedition.}

\textit{The letter is the third or fourth of the Fam. 7
sequence, all of which take their tone from the joke that runs
through them: the urbane Roman man-of-letters lost in the
camp. Trebatius is reported as having despised the
``conveniences of the tribunate'' --- a military tribunate, on
Caesar's staff, which carried pay and standing without combat
duty (``especially with the labour of soldiering taken
away''). Caesar's own report to Cicero --- ``not yet familiar
enough on account of his occupations, but that you certainly
will be'' --- is deflating.}

\textit{The threats of \textsection 2 are mock-juristic:
Vacerra and Manilius are jurists Cicero will go and complain
to; Cornelius is C. Cornelius Maximus, Trebatius's teacher
(``you profess to have learned from him to be wise''), too
weighty to be teased. The closing line --- ``I await your
British letters'' --- is the running theme of the sequence:
Cicero waits, with bemused attention, for the man-of-Roman-law
to send back reports from beyond the world's edge. \textit{Q.
fr.} 2.16 reports Trebatius getting to Britain; \textit{Fam.}
7.10 (a year later) shows him back with Caesar in Gaul,
better-tempered. Trebatius would survive to be one of the
great jurists of the Augustan age and the dedicatee of
\textit{Topica}.}
